%!TEX program = lualatex
\documentclass[ngerman,headheight=70pt]{scrartcl}
\usepackage[ngerman]{babel}
\usepackage{blindtext}
\usepackage[utf8]{luainputenc}
\usepackage{scrlayer-scrpage}
\usepackage[T1]{fontenc}
\usepackage{fontspec}
\usepackage{multicol}
\usepackage{enumitem}
\usepackage[german=quotes]{csquotes}
\usepackage[bottom=2cm,footskip=8mm]{geometry}
\setmainfont{Source Sans Pro}
\MakeOuterQuote{"}

\setlength{\parindent}{0pt}
\setlength{\parskip}{6pt}

\pagestyle{plain.scrheadings}
\thispagestyle{scrheadings}
\lohead{\vspace*{-1.5cm}\includegraphics[scale=0.4]{../images/stupa_logo_rot.pdf}}
\rohead{\\Präsidium des\\Studierendenparlaments}

\lofoot[Universität Hamburg · Präsidium des Studierendenparlaments\\\hrule Von-Melle-Park 5 · D-20146 Hamburg]{Universität Hamburg · Präsidium des Studierendenparlaments\\\hrule Von-Melle-Park 5 · D-20146 Hamburg}
\cofoot[]{}
\cohead[\thepage]{}

\renewcommand{\thesubsection}{TOP \arabic{subsection}}
\renewcommand{\thesubsubsection}{\arabic{subsubsection}.}
\setcounter{subsection}{-1}

\begin{document}
    UHH · StuPa-Präsidium · Von-Melle-Park 5 · D-20146 Hamburg

    %%% PAGE ONE (EDIT LATER) %%%%
    % Hier den Titel des Protokolls eintragen
    \section*{Protokoll der 4. Sitzung des Studierendenparlaments vom 9. Juni 2016}

    % Protokollanten eintragen
    \textbf{Protokoll: Ramon Weilinger}\\
    \textbf{Ort: S08 in der ehemaligen HWP}\\
    \textbf{Beginn: 18.33 Uhr}\\
    \textbf{Ende: 20.43 Uhr}

    % Wohl die komplizierteste Stelle im Protokoll
    % Als Faustregel sollten nicht mehr als 3 Namen auf einer Zeile stehen
    % Einfach an dieser Vorlage orientieren
    % Im Idealfall ist hier nicht viel zu ändern, da kaum jemand fehlt.
    % Rücktritte müssen natürlich pro Protokoll angepasst werden
    \vspace{0.5cm}
    \begin{tabular}{ll}
        Anwesend: & \\
            RCDS (5 Sitze): & Ramon Weilinger, Ramin Shakiba, Jennifer Maack, Benjamin Welling \\
             CampusGrün (14 Sitze): & Laura Franzen, Geoffrey Youett, Elena Rysikova, \\
                                    & Philipp Droll, Yasemin Günther, Melf Johannsen, \\
                                    & Jim Martens, Svenja Horn, Mirzo Ulugbek Khatamov, \\
                                    & Martin Sievert, Armin Günther \\
             BGZM-Liste (2 Sitze) : & Jakob Pape \\
             WiWi (2 Sitze): & \\
             Unicorns (5 Sitze): & Katharina Kucza, Johannes Peplow, Annkathrin Löffler, \\
                                 & Andreas Hartkamp, Marielle Hermstrüwer  \\
             Liste LINKS (3 Sitze): & Gunhild Berdal, Till Petersen, Sinah Mielich \\
             HWP (2 Sitze): &  \\
             MIN (4 Sitze): & Ailina Salten, Jan Detampel \\
             SDS* (3 Sitze): & Mena Winkler, Jacob Petersein, Artur Brückmann \\
             Bart-LISTE (2 Sitze): & Dominic Laumer \\
             LHG (1 Sitz): & Tobias Heisig \\
             harte zeiten (1 Sitz): & Tobias Berking \\
             Jura (1 Sitz): & \\
             AL (2 Sitze): & \\
            & \\
        Entschuldigt: & Antonia Niecke (RCDS), Mario Moldenhauer, Christian Lagod (CampusGrün), \\
                      & Karen Martirosian (BGZM), Claas-Friso Hente (WiWi), Ajdina Karahasan (HWP), \\
                      & Lotte Rullkötter (MIN), Timo Zeimet (Bart), Karima Schulze, \\
                      & Henri Weber (AL)\\
                      &\\
        Unentschuldigt abwesend: & Tahnee Herzig (CampusGrün), Kay Zöllmer (WiWi), Samet Gunay (HWP)\\
                                & Thea Wahlers (MIN), Johann Baumhoefener (Jura)\\
                                &\\
        Rücktritte: & Sandra Fietz (CampusGrün) \rightarrow Christian Lagod\\
    \end{tabular}
    \newpage
    %%% Vorgeschlagene Tagesordnung %%%
    \underline{Vorgeschlagene Tagesordnung}
    \begin{enumerate}[label={\textbf{Top \theenumi}},leftmargin=*]
        % folgende Zeile bei Sitzungsfortsetzungen einkommentieren und x durch
        % die Zahl des zuletzt behandelten TOPs ersetzen
        %\setcounter{enumi}{x}
        \item Keine Werbung von/für Bundeswehr (V1617-023) (60 Min.)
        \item Akkreditierung (V1617-024) (60 Min.)
        \item Prüfauftrag: Wahl-O-Mat bei StuPa-Wahlen (V1617-022) (30 Min.)
        \item Verlegung der StuPa-Wahl in das Sommersemester (V1617-015) (60 Min.)
        \item Verschiedenes
    \end{enumerate}

    \newpage

    % Hier beginnt Teil A. Bei Fortsetzungen kommt hier nur der Teil über
    % die Feststellung der endgültigen Fassung der Tagesordnung.
    % Für sonstige Sitzungen ist dies bereits das Gerüst für die Formalia.
    % Man muss nur noch die protokollierten Inhalte eintragen.

    {\Large \underline{Teil A}}

    \subsection{Formalia}

    \textit{Jim eröffnet die Sitzung um 18.33 Uhr.}

    \subsubsection{Geschäftsbericht Präsidium}

    Jim berichtet aus dem Präsidium:
    \begin{itemize}
        \item Nächste Sitzung des Studierendenparlamentes findet bereits in einer
        Woche am Donnerstag, 16. Juni 2016, statt. Dort wird u.a. der Haushalt in
        erster Lesung beraten.
        \item Sog. "Ersetzungsanträge" werden künftig als eigenständige Vorlagen
        innerhalb eines Tagesordnungspunktes behandelt und dementsprechend
        einzeln nach ihrem Eingang abgestimmt.
    \end{itemize}

    \subsubsection{Anfragen an das Präsidium}

    Keine.

    \subsubsection{Geschäftsbericht AStA}

    Franziska und Philipp berichten für den AStA-Vorstand:
    \begin{itemize}
        \item Kampagne gegen die Exzellenzinitiative: Vollversammlung fand vor
        einer Woche statt mit diskursiver Diskussion. Infolge einer Pressekonferenz
        haben die DPA sowie diverse Medien über die AStA-Haltung berichtet.
        Demonstration gegen die Exzellenzinitiative findet am Dienstag,
        14. Juni 2016, ab 14 Uhr statt. Startpunkt ist vor dem Bahnhof Dammtor.
        \item Es fand ein Treffen mit Vertreter*innen diverser Fachschaftsräte statt,
        um die Begehung durch den Wissenschaftsrat am Montag und Dienstag der
        kommenden Woche vorzubereiten.
        \item Eine Veranstaltung zur Zivilklausel fand statt und war geprägt
        durch eine intensive Diskussion. Das weitere Vorgehen wird nun
        entwickelt werden.
    \end{itemize}

    Timo berichtet aus dem Finanzreferat, dass dieser Tage der Haushaltsplan
    erarbeitet wurde und dass der Haushaltsplanentwurf am Dienstag, 14. Juni 2016,
    ab 18.30 Uhr dem Haushaltsausschuss vorgestellt werde.

    Geoffrey berichtet, dass eine neue Ausgabe des "Stillen Wörtchens" erschienen
    sei. Ferner werde nun die alternative EM-Veranstaltungsreihe realisiert: Immer
    drei Stunden vor Anpfiff des jeweiligen Deutschlandsspiels werde ein Workshop
    stattfinden. Anschließend könne man gemeinsam zum Open-Air-Kino.

    Artur berichtet von den Planungen für eine Veranstaltung zur "Demilitarisierung
    der EU". Eine tiefergehende Debatte hierzu könne im Rahmen der
    Haushaltsberatungen erfolgen.

    Oliver berichtet für den Ausschuss gegen Rechts, dass er allen Fachschaftsräten
    und den Listenverantwortlichen eine E-Mail geschickt habe und warb für
    Reaktionen.

    \subsubsection{Anfragen an den AStA}

    Till fragt nach, ob es keinen Bericht aus dem Referat für Antidiskriminierung gebe.

    Jakob erwidert, dass man in die Vorbereitung der EM-Veranstaltungsreihe
    eingebunden sei.

    Till fragte erneut nach, weshalb es das Referat dann überhaupt als eigenständiges
    Referat gebe, da die Mitarbeit an dieser Veranstaltungsreihe auch ohne
    Referent*innenstatus erfolgen könne.

    Mehrere AstA-Vertreter*innen, darunter Philipp, weisen darauf hin, dass das Referat
    bislang krankheitsbedingt und aufgrund der Ortsabwesenheit von Referent*innen
    nicht mit voller Kraft arbeiten konnte.

    \subsubsection{Dringlichkeitsanträge des AStA}

    Keine.

    \subsubsection{Aktuelle Stunde}

    Es wurde kein Antrag gestellt.

    \subsubsection{Feststellung der endgültigen Fassung des Teils B der Tagesordnung}

    Von Philipp und Till wird vorgeschlagen, die Tagesordnungspunkte 2 (Akkreditierung)
    und 4 (Verlegung StuPa-Wahl) zu vertagen. Dem Vorschlag wird einstimmig bei
    zwei Enthaltungen zugestimmt.

    Die endgültige Fassung des Teils B der Tagesordnung wird einstimmig bei zwei
    Enthaltungen festgestellt.

    %%% BESCHLOSSENE TAGESORDNUNG %%%
    \underline{Endgültige Fassung der Tagesordnung}
    \begin{enumerate}[label={\textbf{Top \theenumi}},leftmargin=*]
        % folgende Zeile bei Sitzungsfortsetzungen einkommentieren und x durch
        % die Zahl des zuletzt behandelten TOPs ersetzen
        %\setcounter{enumi}{x}
        \item Keine Werbung von/für Bundeswehr (V1617-023) (60 Min.)
        \item Prüfauftrag: Wahl-O-Mat bei StuPa-Wahlen (V1617-022) (30 Min.)
        \item Verschiedenes
    \end{enumerate}

    \subsubsection{Feststellung der Beschlussfähigkeit}

    Das Studierendenparlament ist mit 27 anwesenden Mitgliedern beschlussfähig.

    \subsubsection{Genehmigung der Protokolle der vorangegangenen Sitzungen}

    Das Protokoll der dritten Sitzung des Studierendenparlamentes wird mit
    besprochenen Änderungen einstimmig bei einer Enthaltung genehmigt.

    \newpage
    %\vspace{0.5cm}
    {\Large \underline{Teil B}}

    % Teil B beginnt
    % Dieser wird sich für jede Sitzung unterscheiden.
    % Jeder TOP ist eine subsection.
    % Vorhandene Beispiele anfangs entfernen und durch tatsächliche TOPs ersetzen

    % Folgende Zeile bei Sitzungsfortsetzungen einkommentieren und x mit
    % der Zahl des zuletzt behandelten TOPs ersetzen.
    %\setcounter{subsection}{x}
    \subsection{Keine Werbung von/für Bundeswehr}

    Franziska stellt den Antrag für die Antragsteller*innen vor und führt aus, dass der
    Antrag seit der letzten Beratung im Studierendenparlament noch einmal
    überarbeitet wurde.

    Philipp fragt, was mit der Pauschalverurteilung aller Soldat*innen als Mörder
    gemeint ist.

    Gunhild findet es gut, dass der Fokus jetzt mehr auf dem Verhalten der Bundeswehr
    und weniger auf der Jobmesse liege.

    Tobias B. geht insbesondere auf die aktuelle Werbekampagne der Bundeswehr ein
    und verurteilt den Slogan „MACH WAS WIRKLICH ZÄHLT“.

    Till geht auf die Frage von Philipp ein und führt aus, dass die Formulierung
    dahingehend konkretisiert werden sollte, dass Soldat*innen "zum Töten und somit
    als potentiell staatlich sanktionierte Mörder" ausgebildet werden. Ferner
    führte er aus, dass deutlicher werden solle, dass die Bundeswehr kein normaler
    Arbeitsgeber sei und niemand sich gezwungen fühlen sollten, bei der Bundeswehr
    zu arbeiten.

    Artur schlägt vor, dass der AStA beauftragt werden solle, weitergehendes
    Informationsmaterial zu erstellen.

    Till empfiehlt, dass sich der AStA stattdessen intensiv und konzentriert mit
    der Zivilklausel auseinandersetzen sollte.

    \textit{Fraktionspause des RCDS bis 20.05 Uhr.}

    Ein Vertreter aus der Öffentlichkeit, ein Studierender der
    Helmut-Schmidt-Universität, stellt klar, dass es an der Helmut-Schmidt-Universität
    eine studentische Vertretung, nämlich einen studentischen Konvent gebe.

    Till fragt nach, ob es nicht so sei, dass der studentische Konvent nur sehr
    eingeschränkte Möglichkeiten der Meinungskundgabe habe.

    Ein weiterer Student der Helmut-Schmidt-Universität weist darauf hin, dass
    die Studierenden der HSU bereits Angehörige der Bundeswehr seien und daher
    innerhalb dieser Funktion zu politischer Mäßigung gehalten sind.

    Es wird eine Diskussion über die Art und Weise der Erstellung des Antrags
    sowie über die Recherche für die Inhalte des Antrags geführt.

    Alle im Rahmen der Beratungen vorgetragenen Änderungsbegehren wurden von den
    Antragsteller*innen übernommen.

    Der geänderte Gesamtantrag wurde mehrheitlich bei sieben Gegenstimmen und
    drei Enthaltungen beschlossen.

    \textit{Philipp stellt den Antrag TOP 2(neu) "Prüfauftrag: Wahl-O-Mat bei
    StuPa-Wahlen" zu vertagen. Der Antrag wird ohne Widerspruch angenommen.}

    \subsection{Verschiedenes}

    Till weist daraufhin, dass nun verstärkt für die Demonstration gegen die
    Exzellenzinitiative geworben werden müsse. Hierzu seien bereits Plakate bestellt,
    die am Mittag des Folgetages in Hamburg ankämen. Er rief alle Aktiven auf,
    sich am morgigen Freitag, 10. Juni 2016, um 12 Uhr im AStA zu versammeln und
    gemeinsam die Plakate anzubringen. Außerdem empfiehlt er am Montag und
    Dienstag mit Info-Ständen o.Ä. für die Demonstration zu werben.

    Franziska erklärt, dass auch die AStA-Sitzung am Freitag, 10. Juni 2016, um
    14 Uhr direkt mit dem Thema Exzellenzinitiative beginne.

    \textit{Jim schließt um 20.43 Uhr die Sitzung.}

\end{document}
